% ============================================================================
% FLAVOR 3 — FIXPOINT LENS (from fixpoint model checking to model building)
% Entry point: the reader already knows CTL fixpoints; we lift them to dense time
%   and to a model-free, witness-producing setting. Situates the contribution.
% Most "TACAS-positioning" of the three.
% Drop-in replacement for the "Timed modalities" + "Cover rule" subsections.
% Reuses labels tag:Uunf, tag:Runf, tag:chk, tag:chkstar.
% ============================================================================

\subsection{Timed modalities and the cover rule}\label{sec:timed}

\paragraph{The fixpoint view, lifted twice.}
Branching-time model checking rests on fixpoints: $\mathrm{E}(\varphi\,\mathbf{U}\,\psi)$ is a least
fixpoint (reachability of $\psi$), and $\mathrm{AG}\,\psi=\mathrm{A}(\bot\,\mathbf{R}\,\psi)$ a greatest
fixpoint (safety), each evaluated over a fixed structure. Our procedure computes exactly these fixpoints,
lifted in two directions that define the setting. Time is dense, so a fixpoint ranges over \emph{zones}
rather than states. And there is \emph{no} structure, so a fixpoint is computed against the ambient context
$\Ctx$, and its witnesses are what a model is later built from. Concretely, a timed modality
$\mu=\mathrm{Q}(\varphi\,\odot_\kI\,\psi)$ is evaluated to the set
\[
\Sat(\mu)\ \subseteq\ \mathbb{R}_{\ge0}^{\clockSet}
\]
of clock valuations at which it holds under $\Ctx$, defined below as a fixpoint over zone-unions.

\paragraph{The unfolding is the one-step operator.}
Every fixpoint characterisation rests on an unfolding law; ours is the timed, symbolic form of
$\mathrm{E}(\varphi\,\mathbf{U}\,\psi)\equiv \psi\vee(\varphi\wedge \mathrm{EX}\,\mathrm{E}(\varphi\,\mathbf{U}\,\psi))$.
We \emph{pursue} $\mu$ against a formula-clock $x_\mu$ reset at the demand, and one unfolding step rewrites
$\mu^\bullet$ into \emph{base} children, which realise the constant $\psi$-part of the law, and
\emph{recursive} children, which realise the $\mathrm{EX}/\mathrm{AX}$-part by deferring $\nxQ\mu^\bullet$ to
a successor ($\EX\mu^\bullet$ into $\Gamma$ for $\mathrm{Q}=\mathrm{E}$, $\AX\mu^\bullet$ into $\Ctx$ for
$\mathrm{Q}=\mathrm{A}$). For an until, the base child is the stop, and the recursive child is the go, capped
at $x_\mu\le\sup\kI$ in the universal case so that the eventuality cannot escape its window:
\[
\frac{\tup{\Gamma\cup\{\mu^\bullet\}}{\Ctx}{Z}{\obligations}{\eventualities}}
{\begin{array}{ll}
\tup{\Gamma\cup\{\psi\}}{\Ctx}{Z\cap(x_\mu\!\in\!\kI)}{\obligations}{\eventualities\setminus\{\mu\}} & (\text{stop, base})\\[3pt]
\tup{\Gamma\cup\{\varphi,\,\EX\mu^\bullet\}}{\Ctx}{Z}{\obligations}{\eventualities} & (\text{go},\ \mathrm{Q}{=}\mathrm{E})\\[3pt]
\tup{\Gamma\cup\{\varphi\}}{\Ctx\cup\{\AX\mu^\bullet\}}{Z\cap(x_\mu\le\sup\kI)}{\obligations}{\eventualities} & (\text{go-A},\ \mathrm{Q}{=}\mathrm{A})
\end{array}}
\qquad \mu=\mathrm{Q}(\varphi\,\mathbf{U}_\kI\,\psi)
\tag{$\mathbf{U}^\bullet$}\label{tag:Uunf}
\]
For a release the window forces a case split of the zone into before, in, and after $\kI$, with the
discharging cases as base and the continuing cases as recursive:
\[
\frac{\tup{\Gamma\cup\{\mu^\bullet\}}{\Ctx}{Z}{\obligations}{\eventualities}}
{\begin{array}{ll}
\tup{\Gamma\cup\{\varphi\}}{\Ctx}{Z^{<\kI}}{\obligations}{\eventualities} & (\text{before: }\varphi\text{ releases})\\[2pt]
\tup{\Gamma\cup\{\nxQ\mu^\bullet\}}{\Ctx}{Z^{<\kI}}{\obligations}{\eventualities} & (\text{before: continue})\\[2pt]
\tup{\Gamma\cup\{\psi,\,\nxQ\mu^\bullet\}}{\Ctx}{Z^{\in\kI}}{\obligations}{\eventualities} & (\text{in: }\psi,\text{ continue})\\[2pt]
\tup{\Gamma\cup\{\varphi,\psi\}}{\Ctx}{Z^{\in\kI}}{\obligations}{\eventualities} & (\text{in: }\psi,\ \varphi\text{ releases})\\[2pt]
\tup{\Gamma}{\Ctx}{Z^{>\kI}}{\obligations}{\eventualities} & (\text{after: lapsed})
\end{array}}
\qquad \mu=\mathrm{Q}(\varphi\,\mathbf{R}_\kI\,\psi)
\tag{$\mathbf{R}^\bullet$}\label{tag:Runf}
\]
Writing $B_\mu$ for the union of the base-child zones, the recursive children define a monotone operator on
unions of zones,
\[
F_\mu(X)\ =\ B_\mu\ \cup\ \{\,\nu : \text{the successors of $\nu$ required by $\mathrm{Q}$ lie in $X$}\,\},
\]
where ``required'' means \emph{some} successor when $\mathrm{Q}=\mathrm{E}$ and \emph{all} when
$\mathrm{Q}=\mathrm{A}$.

\paragraph{Least for eventualities, greatest for safety.}
$\Sat(\mu)$ is the extremal fixpoint of $F_\mu$: least for an until, whose eventuality must be reached, and
greatest for a release, whose invariant must be maintained,
\[
\Sat(\mu)\ =\
\begin{cases}
\mathrm{lfp}\,F_\mu & \text{if } \odot=\mathbf{U},\\[2pt]
\mathrm{gfp}\,F_\mu & \text{if } \odot=\mathbf{R}.
\end{cases}
\]
Because zones are compared up to $\cong_M$, the lattice of zone-unions is finite, so both fixpoints are
reached by finite Kleene iteration, from $\emptyset$ upward for the least and from the full zone downward
for the greatest, and $\Sat(\mu)$ is a finite union of zones. This is a model-free, context-relative
fixpoint evaluation, and one clock $x_\mu$ absorbs the continuum of onsets of a persistent demand into a
single zone union.

\paragraph{The cover rule: fixpoint membership as a proof step.}
Evaluating $\Sat(\mu)$ answers \emph{where} $\mu$ holds; the cover rule turns that into a step of the
derivation. A demand for $\mu$ on zone $Z$ is met exactly when $Z$ lies inside the fixpoint,
$Z\subseteq\Sat(\mu)$, in which case $\mu$ is discharged and the launched pursuit is retained as the
fixpoint's witness, from which the model reads its runs:
\[
\frac{\tup{\Gamma\cup\{\mu\}}{\Ctx}{Z}{\obligations}{\eventualities}}
{\begin{array}{ll}
\tup{\Gamma}{\Ctx}{Z}{\obligations}{\eventualities} & (\text{discharge})\\[2pt]
\tup{\Gamma\cup\{\mu^\bullet\}}{\Ctx}{Z\cap(x_\mu{=}0)}{\obligations\cup\{x_\mu\}}{\eventualities\cup\{\mu\}} & (\text{launch-U})\\[2pt]
\tup{\Gamma\cup\{\mu^\bullet\}}{\Ctx}{Z\cap(x_\mu{=}0)}{\obligations\cup\{x_\mu\}}{\eventualities} & (\text{launch-R})
\end{array}}
\quad Z\subseteq\Sat(\mu)
\tag{\textrm{Cov}}\label{tag:chk}
\]
On partial cover the demand is refined to $Z\cap\Sat(\mu)$, component by component, and the residual outside
the fixpoint closes:
\[
\frac{\tup{\Gamma\cup\{\mu\}}{\Ctx}{Z}{\obligations}{\eventualities}}
{\tup{\Gamma\cup\{\mu\}}{\Ctx}{Z_i}{\obligations}{\eventualities}}
\qquad Z_i\in\mathsf{cc}\big(Z\cap\Sat(\mu)\big)
\tag{\textrm{Cov}$^{*}$}\label{tag:chkstar}
\]
Where classical model checking tests membership in a fixpoint computed over a \emph{given} structure, here
the identical test is a proof rule whose fixpoint is computed with \emph{no} structure in hand and whose
witness constructs one. That shift, from checking a model to building one, is what the cover rule
contributes.
